\chapter{Parent-потенциал конечного Dirac-оператора}

После факторизации по автоморфизмам конечный оператор сохраняет два модуля
и фазу произведения \(xz\). Минимальный положительный функционал, не
использующий наблюдаемые данные, имеет вид spectral flattening:
\begin{equation}
 V_{\mathrm{flat}}(D_F)
 =\Tr_{\mathrm{orb}}\left(D_F^2-\mathbf1\right)^2.
\end{equation}
На KO6-дублированном модуле orbit trace совпадает со следом исходного
четырёхсекторного квадрата, поэтому не добавляет multiplicity factor.

\section{Точная редукция к орбитным инвариантам}

Положим
\begin{equation}
 r=|x|,
 \qquad
 s=|z|,
 \qquad
 \phi=\arg(xz),
\end{equation}
и введём
\begin{equation}
 u=r^2+s^2,
 \qquad
 v=r^2s^2\sin^2\phi.
\end{equation}
Тогда
\begin{equation}
 V_{\mathrm{flat}}
 =8u^2-16v-8u+4.
\end{equation}
Поскольку исходный функционал является квадратом Hilbert--Schmidt нормы,
\(V_{\mathrm{flat}}\ge0\), и глобальные минимумы удовлетворяют
\begin{equation}
 D_F^2=\mathbf1.
\end{equation}

Для chiral block это эквивалентно
\begin{equation}
 BB^\dagger=\mathbf1_2.
\end{equation}
Диагональные и внедиагональные компоненты дают
\begin{equation}
 2r^2=1,
 \qquad
 2s^2=1,
 \qquad
 2\operatorname{Re}(xz)=0.
\end{equation}
Следовательно,
\begin{equation}
 \boxed{
 r=s=\frac1{\sqrt2},
 \qquad
 \phi=\pm\frac{\pi}{2}\pmod{\pi}.}
\end{equation}

\begin{theorem}[Unitary-flattening orbit selector]
Функционал \(V_{\mathrm{flat}}\) без подбора edge-параметров выбирает
равные модули и максимальную по модулю комплексную фазу. После quotient по
rephasing и обмену листов остаются ровно две CP-сопряжённые орбиты.
\end{theorem}

\section{Hessian-гейт}

В координатах \((r,s,\phi)\) в точке
\begin{equation}
 \left(\frac1{\sqrt2},\frac1{\sqrt2},\frac{\pi}{2}\right)
\end{equation}
гессиан равен
\begin{equation}
 \Hess V_{\mathrm{flat}}
 =
 \begin{pmatrix}
 32&0&0\\
 0&32&0\\
 0&0&8
 \end{pmatrix}.
\end{equation}
Обе радиальные моды и фазовая мода устойчивы. Нулевые направления остаются
только вдоль уже факторизованных gauge/rephasing-орбит.

\section{Сравнение с общим чётным полиномом}

Для
\begin{equation}
 V_{a,b}
 =a\,\Tr_{\mathrm{orb}}D_F^4
 -b\,\Tr_{\mathrm{orb}}D_F^2,
 \qquad a,b>0,
\end{equation}
минимизация также максимизирует \(v\), но масштаб равен
\begin{equation}
 u_{\min}=\frac{b}{2a}.
\end{equation}
Таким образом, равенство модулей и максимальная CP-фаза являются
структурно устойчивыми, а абсолютный радиус орбиты зависит от отношения
коэффициентов. Функционал \(V_{\mathrm{flat}}\) соответствует специальному
отношению \(b/a=2\).

\section{Scale red-team}

Если \(D_F\) имеет физическую размерность массы, корректный функционал
должен содержать масштаб
\begin{equation}
 V_M(D_F)
 =\Tr_{\mathrm{orb}}\left(D_F^2-M^2\mathbf1\right)^2,
\end{equation}
и тогда
\begin{equation}
 |x|=|z|=\frac{M}{\sqrt2}.
\end{equation}
Алгебра, reality и first-order condition не фиксируют \(M\). Запись
\(M=1\) является допустимой только если \(D_F\) заранее определён как
безразмерный нормированный sign/flattened operator. Она не выводит
физический массовый масштаб.

\begin{nogocriterion}[Unit-scale gap]
Unitary-flattening potential считается parameter-free selector только для
безразмерного конечного оператора. Если \(D_F\) является физическим
Dirac/Yukawa-блоком, масштаб \(M\) должен быть выведен из геометрии,
дилатона, cutoff equation или другой части одного parent action.
\end{nogocriterion}

\section{CP-ветвь}

Чётный вещественный функционал зависит от
\(\operatorname{Im}(xz)^2\) и не различает
\(\phi=+\pi/2\) и \(\phi=-\pi/2\). Это точное спонтанное
CP-вырождение. Линейный orientation-odd член мог бы выбрать знак, но его
коэффициент и происхождение должны быть выведены независимо.

\begin{proposition}[Текущий статус]
Parent-потенциал впервые закрывает непрерывную орбиту направления:
отношение модулей и величина CP-фазы выбраны устойчивым минимумом. Полное
замыкание не достигнуто по двум причинам: физический масштаб \(M\) не
выведен, а знак CP-ветви остаётся двукратно вырожденным.
\end{proposition}

\section{Следующий гейт}

Следующий этап должен проверить, может ли уже фиксированная геометрия
\(K\), радиусный сектор или аномальный/eta-вклад:
\begin{enumerate}
 \item сделать \(D_F/M\) канонически безразмерным;
 \item вывести \(M\) без обращения к лептонным массам;
 \item породить orientation-odd вклад либо доказать физическую
 эквивалентность двух CP-ветвей;
 \item одновременно сохранить orbit half-trace.
\end{enumerate}

Машинный аудит сохранён в
\path{s2t_v3_finite_dirac_parent_potential_results.json}.